\IfFileExists{stacks-project.cls}{%
\documentclass{stacks-project}
}{%
\documentclass{amsart}
}

% For dealing with references we use the comment environment
\usepackage{verbatim}
\newenvironment{reference}{\comment}{\endcomment}
%\newenvironment{reference}{}{}
\newenvironment{slogan}{\comment}{\endcomment}
\newenvironment{history}{\comment}{\endcomment}

% For commutative diagrams we use Xy-pic
\usepackage[all]{xy}

% We use 2cell for 2-commutative diagrams.
\xyoption{2cell}
\UseAllTwocells

% We use multicol for the list of chapters between chapters
\usepackage{multicol}

% This is generally recommended for better output
\usepackage{lmodern}
\usepackage[T1]{fontenc}

% For cross-file-references
\usepackage{xr-hyper}

% Package for hypertext links:
\usepackage{hyperref}

% For any local file, say "hello.tex" you want to link to please
% use \externaldocument[hello-]{hello}
\externaldocument[introduction-]{introduction}
\externaldocument[conventions-]{conventions}
\externaldocument[sets-]{sets}
\externaldocument[categories-]{categories}
\externaldocument[topology-]{topology}
\externaldocument[sheaves-]{sheaves}
\externaldocument[sites-]{sites}
\externaldocument[stacks-]{stacks}
\externaldocument[fields-]{fields}
\externaldocument[algebra-]{algebra}
\externaldocument[brauer-]{brauer}
\externaldocument[homology-]{homology}
\externaldocument[derived-]{derived}
\externaldocument[simplicial-]{simplicial}
\externaldocument[more-algebra-]{more-algebra}
\externaldocument[smoothing-]{smoothing}
\externaldocument[modules-]{modules}
\externaldocument[sites-modules-]{sites-modules}
\externaldocument[injectives-]{injectives}
\externaldocument[cohomology-]{cohomology}
\externaldocument[sites-cohomology-]{sites-cohomology}
\externaldocument[dga-]{dga}
\externaldocument[dpa-]{dpa}
\externaldocument[sdga-]{sdga}
\externaldocument[hypercovering-]{hypercovering}
\externaldocument[schemes-]{schemes}
\externaldocument[constructions-]{constructions}
\externaldocument[properties-]{properties}
\externaldocument[morphisms-]{morphisms}
\externaldocument[coherent-]{coherent}
\externaldocument[divisors-]{divisors}
\externaldocument[limits-]{limits}
\externaldocument[varieties-]{varieties}
\externaldocument[topologies-]{topologies}
\externaldocument[descent-]{descent}
\externaldocument[perfect-]{perfect}
\externaldocument[more-morphisms-]{more-morphisms}
\externaldocument[flat-]{flat}
\externaldocument[groupoids-]{groupoids}
\externaldocument[more-groupoids-]{more-groupoids}
\externaldocument[etale-]{etale}
\externaldocument[chow-]{chow}
\externaldocument[intersection-]{intersection}
\externaldocument[pic-]{pic}
\externaldocument[weil-]{weil}
\externaldocument[adequate-]{adequate}
\externaldocument[dualizing-]{dualizing}
\externaldocument[duality-]{duality}
\externaldocument[discriminant-]{discriminant}
\externaldocument[derham-]{derham}
\externaldocument[local-cohomology-]{local-cohomology}
\externaldocument[algebraization-]{algebraization}
\externaldocument[curves-]{curves}
\externaldocument[resolve-]{resolve}
\externaldocument[models-]{models}
\externaldocument[functors-]{functors}
\externaldocument[equiv-]{equiv}
\externaldocument[pione-]{pione}
\externaldocument[etale-cohomology-]{etale-cohomology}
\externaldocument[proetale-]{proetale}
\externaldocument[relative-cycles-]{relative-cycles}
\externaldocument[more-etale-]{more-etale}
\externaldocument[trace-]{trace}
\externaldocument[crystalline-]{crystalline}
\externaldocument[spaces-]{spaces}
\externaldocument[spaces-properties-]{spaces-properties}
\externaldocument[spaces-morphisms-]{spaces-morphisms}
\externaldocument[decent-spaces-]{decent-spaces}
\externaldocument[spaces-cohomology-]{spaces-cohomology}
\externaldocument[spaces-limits-]{spaces-limits}
\externaldocument[spaces-divisors-]{spaces-divisors}
\externaldocument[spaces-over-fields-]{spaces-over-fields}
\externaldocument[spaces-topologies-]{spaces-topologies}
\externaldocument[spaces-descent-]{spaces-descent}
\externaldocument[spaces-perfect-]{spaces-perfect}
\externaldocument[spaces-more-morphisms-]{spaces-more-morphisms}
\externaldocument[spaces-flat-]{spaces-flat}
\externaldocument[spaces-groupoids-]{spaces-groupoids}
\externaldocument[spaces-more-groupoids-]{spaces-more-groupoids}
\externaldocument[bootstrap-]{bootstrap}
\externaldocument[spaces-pushouts-]{spaces-pushouts}
\externaldocument[spaces-chow-]{spaces-chow}
\externaldocument[groupoids-quotients-]{groupoids-quotients}
\externaldocument[spaces-more-cohomology-]{spaces-more-cohomology}
\externaldocument[spaces-simplicial-]{spaces-simplicial}
\externaldocument[spaces-duality-]{spaces-duality}
\externaldocument[formal-spaces-]{formal-spaces}
\externaldocument[restricted-]{restricted}
\externaldocument[spaces-resolve-]{spaces-resolve}
\externaldocument[formal-defos-]{formal-defos}
\externaldocument[defos-]{defos}
\externaldocument[cotangent-]{cotangent}
\externaldocument[examples-defos-]{examples-defos}
\externaldocument[algebraic-]{algebraic}
\externaldocument[examples-stacks-]{examples-stacks}
\externaldocument[stacks-sheaves-]{stacks-sheaves}
\externaldocument[criteria-]{criteria}
\externaldocument[artin-]{artin}
\externaldocument[quot-]{quot}
\externaldocument[stacks-properties-]{stacks-properties}
\externaldocument[stacks-morphisms-]{stacks-morphisms}
\externaldocument[stacks-limits-]{stacks-limits}
\externaldocument[stacks-cohomology-]{stacks-cohomology}
\externaldocument[stacks-perfect-]{stacks-perfect}
\externaldocument[stacks-introduction-]{stacks-introduction}
\externaldocument[stacks-more-morphisms-]{stacks-more-morphisms}
\externaldocument[stacks-geometry-]{stacks-geometry}
\externaldocument[moduli-]{moduli}
\externaldocument[moduli-curves-]{moduli-curves}
\externaldocument[examples-]{examples}
\externaldocument[exercises-]{exercises}
\externaldocument[guide-]{guide}
\externaldocument[desirables-]{desirables}
\externaldocument[coding-]{coding}
\externaldocument[obsolete-]{obsolete}
\externaldocument[fdl-]{fdl}
\externaldocument[index-]{index}

% Theorem environments.
%
\theoremstyle{plain}
\newtheorem{theorem}[subsection]{Theorem}
\newtheorem{proposition}[subsection]{Proposition}
\newtheorem{lemma}[subsection]{Lemma}

\theoremstyle{definition}
\newtheorem{definition}[subsection]{Definition}
\newtheorem{example}[subsection]{Example}
\newtheorem{exercise}[subsection]{Exercise}
\newtheorem{situation}[subsection]{Situation}

\theoremstyle{remark}
\newtheorem{remark}[subsection]{Remark}
\newtheorem{remarks}[subsection]{Remarks}

\numberwithin{equation}{subsection}

% Macros
%
\def\lim{\mathop{\mathrm{lim}}\nolimits}
\def\colim{\mathop{\mathrm{colim}}\nolimits}
\def\Spec{\mathop{\mathrm{Spec}}}
\def\Hom{\mathop{\mathrm{Hom}}\nolimits}
\def\Ext{\mathop{\mathrm{Ext}}\nolimits}
\def\SheafHom{\mathop{\mathcal{H}\!\mathit{om}}\nolimits}
\def\SheafExt{\mathop{\mathcal{E}\!\mathit{xt}}\nolimits}
\def\Sch{\mathit{Sch}}
\def\Mor{\mathop{\mathrm{Mor}}\nolimits}
\def\Ob{\mathop{\mathrm{Ob}}\nolimits}
\def\Sh{\mathop{\mathit{Sh}}\nolimits}
\def\NL{\mathop{N\!L}\nolimits}
\def\CH{\mathop{\mathrm{CH}}\nolimits}
\def\proetale{{pro\text{-}\acute{e}tale}}
\def\etale{{\acute{e}tale}}
\def\QCoh{\mathit{QCoh}}
\def\Ker{\mathop{\mathrm{Ker}}}
\def\Im{\mathop{\mathrm{Im}}}
\def\Coker{\mathop{\mathrm{Coker}}}
\def\Coim{\mathop{\mathrm{Coim}}}

% Boxtimes
%
\DeclareMathSymbol{\boxtimes}{\mathbin}{AMSa}{"02}

%
% Macros for moduli stacks/spaces
%
\def\QCohstack{\mathcal{QC}\!\mathit{oh}}
\def\Cohstack{\mathcal{C}\!\mathit{oh}}
\def\Spacesstack{\mathcal{S}\!\mathit{paces}}
\def\Quotfunctor{\mathrm{Quot}}
\def\Hilbfunctor{\mathrm{Hilb}}
\def\Curvesstack{\mathcal{C}\!\mathit{urves}}
\def\Polarizedstack{\mathcal{P}\!\mathit{olarized}}
\def\Complexesstack{\mathcal{C}\!\mathit{omplexes}}
% \Pic is the operator that assigns to X its picard group, usage \Pic(X)
% \Picardstack_{X/B} denotes the Picard stack of X over B
% \Picardfunctor_{X/B} denotes the Picard functor of X over B
\def\Pic{\mathop{\mathrm{Pic}}\nolimits}
\def\Picardstack{\mathcal{P}\!\mathit{ic}}
\def\Picardfunctor{\mathrm{Pic}}
\def\Deformationcategory{\mathcal{D}\!\mathit{ef}}
\usepackage{cleveref}
\crefname{originalreference}{original reference}{original references}
\Crefname{originalreference}{Original reference}{Original references}
\makeatletter
\newlabel{lemma-simplicial-cocontinuous-site}{{09WD}{source}{}{sourcepointer-1}{}}
\newlabel{lemma-simplicial-cocontinuous-site@cref}{{[originalreference][0][]09WD}{[1][0][]source}{}{}{}}
\makeatother
\crefname{originalreference}{original reference}{original references}
\Crefname{originalreference}{Original reference}{Original references}
\makeatletter
\newlabel{lemma-simplicial-site-site}{{09WC}{source}{}{sourcepointer-2}{}}
\newlabel{lemma-simplicial-site-site@cref}{{[originalreference][0][]09WC}{[1][0][]source}{}{}{}}
\makeatother
\crefname{originallemma}{original lemma}{original lemmas}
\Crefname{originallemma}{Original lemma}{Original lemmas}
\makeatletter
\newlabel{sites-lemma-cocontinuous-morphism-topoi}{{00XO}{source}{}{sourcepointer-3}{}}
\newlabel{sites-lemma-cocontinuous-morphism-topoi@cref}{{[originallemma][0][]00XO}{[1][0][]source}{}{}{}}
\makeatother
\crefname{originallemma}{original lemma}{original lemmas}
\Crefname{originallemma}{Original lemma}{Original lemmas}
\makeatletter
\newlabel{sites-lemma-when-shriek}{{00XR}{source}{}{sourcepointer-4}{}}
\newlabel{sites-lemma-when-shriek@cref}{{[originallemma][0][]00XR}{[1][0][]source}{}{}{}}
\makeatother
\crefname{originallemma}{original lemma}{original lemmas}
\Crefname{originallemma}{Original lemma}{Original lemmas}
\makeatletter
\newlabel{sites-modules-lemma-g-shriek-adjoint}{{04BG}{source}{}{sourcepointer-5}{}}
\newlabel{sites-modules-lemma-g-shriek-adjoint@cref}{{[originallemma][0][]04BG}{[1][0][]source}{}{}{}}
\makeatother
\crefname{originalreference}{original reference}{original references}
\Crefname{originalreference}{Original reference}{Original references}
\makeatletter
\newlabel{situation-simplicial-site}{{09WE}{source}{}{sourcepointer-6}{}}
\newlabel{situation-simplicial-site@cref}{{[originalreference][0][]09WE}{[1][0][]source}{}{}{}}
\makeatother
\title{Component inclusions into the simplicial total site are continuous and cocontinuous, with explicit shriek functors}
\author{The Stacks Project Authors}
\date{stacks-project:spaces-simplicial.tex:local-snapshot-20261012; GNU FDL 1.2 or later}
\begin{document}
\maketitle
\noindent\textit{Editorial note.} All delivered mathematical text is the author's own, copied from the pinned Stacks Project chapter file sources/stacks-spaces-simplicial/pinned-original-source.tex; no statement or proof step was rewritten, reordered or removed.\par\medskip
\noindent\textit{Editorial note.} The counted proof is the author's complete argument: the single omission (source line 1056, 'we omit the arguments showing these two constructions give mutually inverse maps') is the inverse check of the adjunction bijection, auxiliary to parts (3) and (6) which are constructed in full; the author's typo 'ammner' (source line 1066, for 'manner') is preserved verbatim; the left adjoints g\_\{n!\} established here, with their component formula, are used by card L46148 of this batch.\par\medskip
\noindent\textit{Editorial note.} The label lemma-restriction-to-components-site names the counted statement, and the proof environment immediately following it in the source is delivered verbatim as its proof.\par\medskip
\noindent\textit{Editorial note.} The counted statement is carried by the author's own lemma environment in the chapter source.\par\medskip
\noindent\textit{Allowed standard prerequisites (editorial).}\begin{itemize}
\item Morphisms of topoi from cocontinuous functors. Sites chapter, \textbackslash{}label\{sites-lemma-cocontinuous-morphism-topoi\}. Line 1045: produces the morphism of topoi g\_n and part (1) in both cases A and B.
\item Existence of the shriek functor for a cocontinuous functor. Sites chapter, \textbackslash{}label\{sites-lemma-when-shriek\}. Line 1046: existence of the sheaf-level left adjoint g\textasciicircum{}\{Sh\}\_\{n!\} (part (2)).
\item The shriek adjoint on abelian sheaves. Modules on Sites chapter, \textbackslash{}label\{sites-modules-lemma-g-shriek-adjoint\}. Line 1048: existence of the abelian left adjoint g\_\{n!\} (part (5)).
\item Coverings and fibre products in the simplicial total site. Same chapter, \textbackslash{}label\{lemma-simplicial-site-site\} and \textbackslash{}label\{lemma-simplicial-cocontinuous-site\}. Lines 995-1041: in case A resp. case B, images of coverings are coverings by definition and fibre products computed in C\_n agree with those in C\_total (the claim in the cited lemma's proof).
\end{itemize}\medskip
\ifdefined\hypertarget\hypertarget{sourcepointer-1}{}\fi
\noindent\textit{Source reference.} Original reference 09WD: sources/stacks-spaces-simplicial/pinned-original-source.tex, source lines 844--844. Prerequisite pointer invoked at proof line(s) 1009; it names a standard fact of the same chapter that the counted argument uses, and every step of the counted proof built on it is retained verbatim in the delivered body.\par\medskip
\ifdefined\hypertarget\hypertarget{sourcepointer-2}{}\fi
\noindent\textit{Source reference.} Original reference 09WC: sources/stacks-spaces-simplicial/pinned-original-source.tex, source lines 797--797. Prerequisite pointer invoked at proof line(s) 996; it names a standard fact of the same chapter that the counted argument uses, and every step of the counted proof built on it is retained verbatim in the delivered body.\par\medskip
\ifdefined\hypertarget\hypertarget{sourcepointer-3}{}\fi
\noindent\textit{Source reference.} Original lemma 00XO: sources/stacks-chow/tags.txt, source lines 1226--1226. Prerequisite pointer invoked at proof line(s) 1022; it names a standard fact of a sibling chapter that the counted argument uses, and every step of the counted proof built on it is retained verbatim in the delivered body.\par\medskip
\ifdefined\hypertarget\hypertarget{sourcepointer-4}{}\fi
\noindent\textit{Source reference.} Original lemma 00XR: sources/stacks-chow/tags.txt, source lines 1229--1229. Prerequisite pointer invoked at proof line(s) 1023; it names a standard fact of a sibling chapter that the counted argument uses, and every step of the counted proof built on it is retained verbatim in the delivered body.\par\medskip
\ifdefined\hypertarget\hypertarget{sourcepointer-5}{}\fi
\noindent\textit{Source reference.} Original lemma 04BG: sources/stacks-chow/tags.txt, source lines 5651--5651. Prerequisite pointer invoked at proof line(s) 1025; it names a standard fact of a sibling chapter that the counted argument uses, and every step of the counted proof built on it is retained verbatim in the delivered body.\par\medskip
\ifdefined\hypertarget\hypertarget{sourcepointer-6}{}\fi
\noindent\textit{Source reference.} Original reference 09WE: sources/stacks-spaces-simplicial/pinned-original-source.tex, source lines 888--888. Prerequisite pointer invoked at proof line(s) 971; it names a standard fact of the same chapter that the counted argument uses, and every step of the counted proof built on it is retained verbatim in the delivered body.\par\medskip
\begin{lemma}
\label{lemma-restriction-to-components-site}
In Situation \ref{situation-simplicial-site} the functor
$\mathcal{C}_n \to \mathcal{C}_{total}$, $U \mapsto U$ is continuous
and cocontinuous. The associated morphism of
topoi $g_n : \Sh(\mathcal{C}_n) \to \Sh(\mathcal{C}_{total})$ satisfies
\begin{enumerate}
\item $g_n^{-1}$ associates to the sheaf $\mathcal{F}$ on $\mathcal{C}_{total}$
the sheaf $\mathcal{F}_n$ on $\mathcal{C}_n$,
\item $g_n^{-1} : \Sh(\mathcal{C}_{total}) \to \Sh(\mathcal{C}_n)$
has a left adjoint $g^{Sh}_{n!}$,
\item for $\mathcal{G}$ in $\Sh(\mathcal{C}_n)$ the restriction of
$g_{n!}^{Sh}\mathcal{G}$ to $\mathcal{C}_m$ is
$\coprod\nolimits_{\varphi : [n] \to [m]} f_\varphi^{-1}\mathcal{G}$,
\item $g_{n!}^{Sh}$ commutes with finite connected limits,
\item $g_n^{-1} : \textit{Ab}(\mathcal{C}_{total}) \to
\textit{Ab}(\mathcal{C}_n)$ has a left adjoint $g_{n!}$,
\item for $\mathcal{G}$ in $\textit{Ab}(\mathcal{C}_n)$ the restriction of
$g_{n!}\mathcal{G}$ to $\mathcal{C}_m$ is
$\bigoplus\nolimits_{\varphi : [n] \to [m]} f_\varphi^{-1}\mathcal{G}$, and
\item $g_{n!}$ is exact.
\end{enumerate}
\end{lemma}

\begin{proof}
Case A. If $\{U_i \to U\}_{i \in I}$ is a covering in $\mathcal{C}_n$
then the image $\{U_i \to U\}_{i \in I}$ is a covering in $\mathcal{C}_{total}$
by definition (Lemma \ref{lemma-simplicial-site-site}). For a morphism
$V \to U$ of $\mathcal{C}_n$, the fibre product
$V \times_U U_i$ in $\mathcal{C}_n$ is also
the fibre product in $\mathcal{C}_{total}$ (by the claim in the
proof of Lemma \ref{lemma-simplicial-site-site}).
Therefore our functor is continuous. On the other hand, our functor
defines a bijection between coverings of $U$ in $\mathcal{C}_n$
and coverings of $U$ in $\mathcal{C}_{total}$. Therefore it is
certainly the case that our functor is cocontinuous.

\medskip\noindent
Case B. If $\{U_i \to U\}_{i \in I}$ is a covering in $\mathcal{C}_n$
then the image $\{U_i \to U\}_{i \in I}$ is a covering in $\mathcal{C}_{total}$
by definition (Lemma \ref{lemma-simplicial-cocontinuous-site}). For a morphism
$V \to U$ of $\mathcal{C}_n$, the fibre product
$V \times_U U_i$ in $\mathcal{C}_n$ is also
the fibre product in $\mathcal{C}_{total}$ (by the claim in the
proof of Lemma \ref{lemma-simplicial-cocontinuous-site}).
Therefore our functor is continuous. On the other hand, our functor
defines a bijection between coverings of $U$ in $\mathcal{C}_n$
and coverings of $U$ in $\mathcal{C}_{total}$. Therefore it is
certainly the case that our functor is cocontinuous.

\medskip\noindent
At this point part (1) and the existence of $g^{Sh}_{n!}$ and $g_{n!}$
in cases A and B follows from
Sites, Lemmas \ref{sites-lemma-cocontinuous-morphism-topoi} and
\ref{sites-lemma-when-shriek}
and
Modules on Sites, Lemma \ref{sites-modules-lemma-g-shriek-adjoint}.

\medskip\noindent
Proof of (3). Let $\mathcal{G}$ be a sheaf on $\mathcal{C}_n$.
Consider the sheaf $\mathcal{H}$ on $\mathcal{C}_{total}$
whose degree $m$ part is the sheaf
$$
\mathcal{H}_m = \coprod\nolimits_{\varphi : [n] \to [m]}
f_\varphi^{-1}\mathcal{G}
$$
given in part (3) of the statement of the lemma.
Given a map $\psi : [m] \to [m']$ the map
$\mathcal{H}(\psi) : f_\psi^{-1}\mathcal{H}_m \to \mathcal{H}_{m'}$
is given on components by the identifications
$$
f_\psi^{-1} f_\varphi^{-1} \mathcal{G} \to
f_{\psi \circ \varphi}^{-1}\mathcal{G}
$$
Observe that given a map $\alpha : \mathcal{H} \to \mathcal{F}$
of sheaves on $\mathcal{C}_{total}$ we obtain a map
$\mathcal{G} \to \mathcal{F}_n$
corresponding to the restriction of $\alpha_n$ to the component
$\mathcal{G}$ in $\mathcal{H}_n$. Conversely, given a map
$\beta : \mathcal{G} \to \mathcal{F}_n$ of sheaves on $\mathcal{C}_n$
we can define
$\alpha : \mathcal{H} \to \mathcal{F}$ by letting $\alpha_m$
be the map which on components
$$
f_\varphi^{-1}\mathcal{G} \to \mathcal{F}_m
$$
uses the maps adjoint to $\mathcal{F}(\varphi) \circ f_\varphi^{-1}\beta$.
We omit the arguments showing these two constructions give
mutually inverse maps
$$
\Mor_{\Sh(\mathcal{C}_n)}(\mathcal{G}, \mathcal{F}_n) =
\Mor_{\Sh(\mathcal{C}_{total})}(\mathcal{H}, \mathcal{F})
$$
Thus $\mathcal{H} = g^{Sh}_{n!}\mathcal{G}$ as desired.

\medskip\noindent
Proof of (4). If $\mathcal{G}$ is an abelian sheaf on $\mathcal{C}_n$,
then we proceed in exactly the same ammner as above, except that
we define $\mathcal{H}$ is the abelian sheaf on $\mathcal{C}_{total}$
whose degree $m$ part is the sheaf
$$
\bigoplus\nolimits_{\varphi : [n] \to [m]} f_\varphi^{-1}\mathcal{G}
$$
with transition maps defined exactly as above. The bijection
$$
\Mor_{\textit{Ab}(\mathcal{C}_n)}(\mathcal{G}, \mathcal{F}_n) =
\Mor_{\textit{Ab}(\mathcal{C}_{total})}(\mathcal{H}, \mathcal{F})
$$
is proved exactly as above.
Thus $\mathcal{H} = g_{n!}\mathcal{G}$ as desired.

\medskip\noindent
The exactness properties of $g^{Sh}_{n!}$ and $g_{n!}$ follow
from formulas given for these functors.
\end{proof}

\end{document}
